\documentclass[11pt]{article}
\usepackage[margin=1in]{geometry}
\usepackage{amsmath,amssymb,amsthm,mathtools,mathrsfs}
\usepackage{graphicx}
\usepackage{microtype}
\usepackage{enumitem}
\usepackage{xcolor}
\usepackage[colorlinks=true,linkcolor=blue!55!black,citecolor=blue!55!black,urlcolor=blue!65!black]{hyperref}

\newtheorem{theorem}{Theorem}
\newtheorem{remark}{Remark}
\newcommand{\ellp}{\ell_2}
\newcommand{\K}{\mathscr K}

\title{Sublinear Kalton--Peck centralizers are not singular}
\author{Candidate negative answer to the question in arXiv:1410.5505}
\date{19 July 2026}

\begin{document}
\maketitle

\noindent\fbox{\parbox{0.94\textwidth}{
\textbf{Status: candidate full negative answer; human review required.}
For every $0<r<1$, the Kalton--Peck centralizer $\K_{\phi_r}$ on $\ell_2$ is
not singular.  The proof exhibits a disjointly generated copy of $\ell_2$ on
which $\K_{\phi_r}$ is trivial.
}}

\section{Source and exact question}

The source is J.~M.~F. Castillo, V. Ferenczi, and M. Gonz\'alez,
\emph{Singular twisted sums generated by complex interpolation},
arXiv:1410.5505v1 (2014) \cite{source}.  At the end of Section 6, page 25 of
the arXiv PDF, the authors ask whether the sublinear Kalton--Peck maps are
singular:

\begin{figure}[h]
  \centering
  \includegraphics[width=\textwidth]{figures/open_problem_crop.png}
  \caption{The open question rendered from page 25 of the source PDF.}
\end{figure}

\noindent\textbf{Transcription.}
``We already know [32, Corollary 5.5] that $\mathscr K=\mathscr K_{\phi_1}$
is singular but, since the function $\phi_r$ is not expansive for $r<1$, we
do not know if also $\mathscr K_{\phi_r}$ is singular for $0<r<1$.''

\section{Definitions}

For $0<r\leq1$, set
\[
 \phi_r(t)=
 \begin{cases}
 t,&0\leq t\leq1,\\
 t^r,&t>1.
 \end{cases}
\]
The associated Kalton--Peck centralizer on $\ell_2$ is
\[
 \K_{\phi_r}(x)
 =x\,\phi_r\!\left(-\log\frac{|x|}{\|x\|_2}\right),
\]
with all operations taken coordinatewise.  A quasi-linear map is
\emph{trivial on a subspace $Y$} if there is a linear map $L:Y\to\mathbb C^{\mathbb N}$
such that $\K_{\phi_r}-L:Y\to\ell_2$ is bounded.  It is singular exactly when
its restriction to every infinite-dimensional closed subspace is nontrivial;
equivalently, the quotient map in the associated twisted sum is strictly
singular \cite[Section 4]{source}.

\section{Proof Intuition}

On a normalized flat block of length $m$, every nonzero coordinate has
logarithmic size $s=\frac12\log m$.  A linear diagonal map can remove the base
value $s^r$ from the centralizer.  For a linear combination with coefficient
$a$, the remaining nonlinear error is
\[
 (s-\log|a|)^r-s^r.
\]
Because $r<1$, increments of $t^r$ flatten as the base point $s$ grows.  We
choose successive block sizes so rapidly increasing that these increments are
square summable uniformly over every unit vector.  The only apparent danger is
from extremely small coefficients, but their exponential weight
$|a|^2=e^{-2t}$ beats the power $t^{2r}$.  This produces an explicit
infinite-dimensional subspace on which the centralizer is trivial.

\section{Negative answer}

\begin{theorem}\label{thm:main}
For every $0<r<1$, the centralizer $\K_{\phi_r}$ on $\ell_2$ is not singular.
More precisely, there are pairwise disjoint finite sets $A_n\subset\mathbb N$
such that, for
\[
 u_n=|A_n|^{-1/2}\mathbf 1_{A_n},\qquad
 Y=\overline{\operatorname{span}}\{u_n:n\geq1\},
\]
the restriction $\K_{\phi_r}|_Y$ is trivial.  Consequently, the quotient map
of $\ell_2(\phi_r)=\ell_2\oplus_{\K_{\phi_r}}\ell_2$ is not strictly singular.
\end{theorem}

\begin{proof}
Fix $0<r<1$ and put
\[
 \beta=2(1-r)>0.
\]
Choose $a>1/(2\beta)=1/(4(1-r))$.  Take pairwise disjoint finite sets $A_n$
with cardinalities $m_n$ so large that
\[
 s_n:=\frac12\log m_n\geq\max\{1,n^a\}.
\]
The disjoint normalized vectors $u_n=m_n^{-1/2}\mathbf1_{A_n}$ are isometric
to the unit vector basis of $\ell_2$, so $Y\cong\ell_2$.

Define a linear, coordinatewise map $L:Y\to\mathbb C^{\mathbb N}$ by
\[
 L(u_n)=s_n^r u_n.
\]
We prove that $\K_{\phi_r}-L$ is bounded from $Y$ to $\ell_2$.
By homogeneity it suffices to consider
\[
 x=\sum_n a_nu_n,\qquad \sum_n|a_n|^2=1.
\]
For $a_n\neq0$, write
\[
 p_n=|a_n|^2,\qquad
 t_n=-\log|a_n|=\frac12\log\frac1{p_n}.
\]
Each coordinate of $x$ in $A_n$ has magnitude $|a_n|m_n^{-1/2}$, hence its
negative logarithm is $s_n+t_n\geq1$.  Both arguments therefore lie in the
power branch of $\phi_r$, and disjointness gives the exact identity
\begin{equation}\label{eq:error}
 \|(\K_{\phi_r}-L)x\|_2^2
 =\sum_{n:p_n>0}p_n\big((s_n+t_n)^r-s_n^r\big)^2.
\end{equation}

Split the indices into
\[
 I=\{n:t_n\leq s_n\},\qquad J=\{n:t_n>s_n\}.
\]
For $n\in I$, concavity and the mean value theorem yield
\[
 (s_n+t_n)^r-s_n^r\leq r t_ns_n^{r-1}.
\]
There is an absolute $C$ such that
\[
 p\log^2(1/p)\leq C\sqrt p\qquad(0<p\leq1).
\]
Since $t_n=\frac12\log(1/p_n)$, Cauchy--Schwarz gives
\begin{align*}
 \sum_{n\in I}p_n\big((s_n+t_n)^r-s_n^r\big)^2
 &\leq r^2\sum_n p_nt_n^2s_n^{-\beta}\\
 &\leq C_r\sum_n s_n^{-\beta}\sqrt{p_n}\\
 &\leq C_r\Big(\sum_n s_n^{-2\beta}\Big)^{1/2}
             \Big(\sum_n p_n\Big)^{1/2}<\infty.
\end{align*}
The final series converges because $s_n\geq n^a$ and $2a\beta>1$.

For $n\in J$, subadditivity of $t\mapsto t^r$ gives
\[
 (s_n+t_n)^r-s_n^r\leq t_n^r.
\]
Also $p_n=e^{-2t_n}$.  The function $e^{-2t}t^{2r}$ decreases for $t\geq r$;
as $t_n>s_n\geq1>r$, we obtain
\[
 p_nt_n^{2r}=e^{-2t_n}t_n^{2r}
 \leq e^{-2s_n}s_n^{2r}.
\]
Consequently,
\[
 \sum_{n\in J}p_n\big((s_n+t_n)^r-s_n^r\big)^2
 \leq\sum_n e^{-2s_n}s_n^{2r}<\infty.
\]
Both bounds are independent of the normalized vector $x$.  Equation
\eqref{eq:error} therefore proves
\[
 \|\K_{\phi_r}(x)-Lx\|_2\leq C_r\|x\|_2\qquad(x\in Y),
\]
so $\K_{\phi_r}|_Y$ is trivial.

For completeness, let $q:\ell_2\oplus_{\K_{\phi_r}}\ell_2\to\ell_2$ be the
canonical quotient and define $T:Y\to\ell_2(\phi_r)$ by $T(x)=(Lx,x)$.
Then $qT=\operatorname{id}_Y$, while
\[
 \|T(x)\|=\|Lx-\K_{\phi_r}(x)\|_2+\|x\|_2
 \asymp\|x\|_2.
\]
Thus $T(Y)$ is closed and $q|_{T(Y)}$ is an isomorphism onto $Y$.  Hence $q$
is not strictly singular.
\end{proof}

\begin{remark}
The endpoint $r=1$ is excluded for a structural reason: then $\beta=0$, the
summability mechanism disappears, and the classical Kalton--Peck centralizer
is singular, as recalled in the source.
\end{remark}

\section{Verification and scope}

The accompanying \texttt{verification.md} separately audits the block formula,
both summability regions, the role of the ambient-valued linear map, and the
passage from triviality to failure of strict singularity.  No computational
experiment is used as proof.

The result answers exactly the source's question for every $0<r<1$.  The block
subspace depends on $r$, and no uniform subspace over the whole interval is
claimed.

\section{Bounded novelty check and review recommendation}

On 19 July 2026, the run's registry, solution, attempt, and proof-gap indexes
were searched for arXiv:1410.5505, $\ell_2(\phi_r)$, sublinear Kalton--Peck
centralizers, and the exact singularity phrase.  Exact-phrase and title/author
web searches were also run.  The most relevant later primary result located,
arXiv:2408.07827, concerns singular spaces built from bi-Lipschitz functions;
it does not answer the sublinear $t^r$ question.  No matching answer was found.
This bounded search does not establish originality.

\medskip
\noindent\textbf{Human-review recommendation.}
Treat this as a high-priority candidate full negative answer.  The decisive
checks are the convention allowing $L:Y\to\mathbb C^{\mathbb N}$ in triviality,
the exponent $a>1/(4(1-r))$, and the explicit lifted copy $T(Y)$ in the twisted
sum.

\begin{thebibliography}{9}
\bibitem{source}
J.~M.~F. Castillo, V. Ferenczi, and M. Gonz\'alez,
\emph{Singular twisted sums generated by complex interpolation},
arXiv:1410.5505v1 (2014),
\url{https://arxiv.org/abs/1410.5505}.
\end{thebibliography}

\end{document}
